\section{Двумерные случайные величины и корреляция}

\subsection{Совместное распределение и математическое ожидание суммы}

Многие эксперименты требуют изучения сразу двух параметров $X$ и $Y$. Совместное распределение пары дискретных случайных величин задается матрицей вероятностей $p_{ij} = P(X = x_i, Y = y_j)$.

\begin{theorem}[Линейность математического ожидания для суммы]
Для любых двух случайных величин $X$ и $Y$ (независимо от того, зависимы они или нет) выполняется:
$$E(X + Y) = E(X) + E(Y).$$
\end{theorem}

\begin{proof}
По определению математического ожидания функции от двумерной случайной величины, мы должны просуммировать значения суммы $(x_i + y_j)$ по всем возможным исходам с их вероятностями $p_{ij}$:
\begin{align*}
E(X + Y) &= \sum_{i} \sum_{j} (x_i + y_j) p_{ij} = \sum_{i} \sum_{j} (x_i p_{ij} + y_j p_{ij}) \\
&= \sum_{i} \sum_{j} x_i p_{ij} + \sum_{i} \sum_{j} y_j p_{ij}.
\end{align*}
В первой двойной сумме вынесем $x_i$ (оно не зависит от индекса $j$), а во второй вынесем $y_j$ (не зависит от $i$):
$$E(X + Y) = \sum_{i} x_i \left( \sum_{j} p_{ij} \right) + \sum_{j} y_j \left( \sum_{i} p_{ij} \right).$$
Заметим, что суммирование совместной вероятности $p_{ij}$ по всем $j$ дает маргинальную (одномерную) вероятность $P(X = x_i)$. Аналогично, сумма по $i$ дает $P(Y = y_j)$:
$$E(X + Y) = \sum_{i} x_i P(X = x_i) + \sum_{j} y_j P(Y = y_j) = E(X) + E(Y).$$
Теорема доказана.
\end{proof}

\begin{theorem}[Математическое ожидание произведения независимых СВ]
Если $X$ и $Y$ независимы, то $E(XY) = E(X)E(Y)$.
\end{theorem}

\begin{proof}
Если $X$ и $Y$ независимы, то вероятность совместного наступления равна произведению вероятностей: $p_{ij} = P(X = x_i) \cdot P(Y = y_j)$.
$$E(XY) = \sum_i \sum_j (x_i y_j) p_{ij} = \sum_i \sum_j (x_i y_j) P(X = x_i) P(Y = y_j).$$
Сгруппируем члены, зависящие от $i$ и от $j$:
$$E(XY) = \left( \sum_i x_i P(X = x_i) \right) \cdot \left( \sum_j y_j P(Y = y_j) \right) = E(X) \cdot E(Y).$$
\end{proof}

\subsection{Ковариация и корреляция}

\begin{definition}
\textbf{Ковариацией} случайных величин $X$ и $Y$ называется математическое ожидание произведения их отклонений от своих математических ожиданий:
$$\Cov(X, Y) = E\left((X - E(X))(Y - E(Y))\right).$$
\end{definition}

\begin{theorem}[Вычислительная формула и независимость]
\begin{enumerate}
\item $\Cov(X, Y) = E(XY) - E(X)E(Y)$.
\item Если $X$ и $Y$ независимы, то $\Cov(X, Y) = 0$.
\end{enumerate}
\end{theorem}

\begin{proof}
\begin{enumerate}
\item Пусть $\mu_X = E(X)$ и $\mu_Y = E(Y)$. Раскроем скобки:
\begin{align*}
\Cov(X,Y) &= E((X - \mu_X)(Y - \mu_Y)) = E(XY - X\mu_Y - Y\mu_X + \mu_X\mu_Y) \\
&= E(XY) - \mu_Y E(X) - \mu_X E(Y) + \mu_X\mu_Y \\
&= E(XY) - \mu_Y\mu_X - \mu_X\mu_Y + \mu_X\mu_Y = E(XY) - \mu_X\mu_Y = E(XY) - E(X)E(Y).
\end{align*}

\item Если $X$ и $Y$ независимы, то по доказанной выше теореме $E(XY) = E(X)E(Y)$.
Подставляем это в формулу: $\Cov(X,Y) = E(X)E(Y) - E(X)E(Y) = 0$.
\end{enumerate}
\end{proof}

\begin{theorem}[Дисперсия суммы]
$D(X + Y) = D(X) + D(Y) + 2\,\Cov(X, Y)$.
\end{theorem}

\begin{proof}
По определению дисперсии: $D(X+Y) = E\left( ((X+Y) - E(X+Y))^2 \right)$.
Перегруппируем слагаемые:
$$D(X+Y) = E\left( ((X - E(X)) + (Y - E(Y)))^2 \right).$$
Возведем в квадрат по формуле $(a+b)^2 = a^2 + 2ab + b^2$:
\begin{align*}
D(X+Y) &= E\left( (X - E(X))^2 + 2(X - E(X))(Y - E(Y)) + (Y - E(Y))^2 \right) \\
&= E\left((X - E(X))^2\right) + 2E\left((X - E(X))(Y - E(Y))\right) + E\left((Y - E(Y))^2\right) \\
&= D(X) + 2\Cov(X, Y) + D(Y).
\end{align*}
\end{proof}